% =====================================================================
% CSE 341: Concepts of Programming Languages
% Project D1 (Part 1): Design Specification for ANS
% Student: 220104004931 Enes Özkan
% Compile with: tectonic 220104004931_D1_P1.tex
% =====================================================================
\documentclass[11pt,a4paper]{article}

\usepackage[utf8]{inputenc}
\usepackage[T1]{fontenc}
\usepackage{mathpazo}
\usepackage[scaled=0.92]{helvet}
\usepackage[varqu,varl]{zi4}
\usepackage[english]{babel}
\usepackage{microtype}
\usepackage[a4paper,margin=2.25cm,top=2.0cm,bottom=2.15cm]{geometry}

\usepackage[table]{xcolor}
\definecolor{ansblue}{HTML}{1E3A8A}
\definecolor{anscodebg}{HTML}{F9FAFB}
\definecolor{anscaption}{HTML}{6B7280}
\definecolor{anskeyword}{HTML}{2563EB}
\usepackage[colorlinks=true, linkcolor=ansblue, urlcolor=ansblue,
            citecolor=ansblue, pdfborder={0 0 0}]{hyperref}

\usepackage[explicit]{titlesec}
\titleformat{\section}
  {\large\bfseries\sffamily\color{black}}
  {\thesection}{0.65em}{#1}
\titleformat{\subsection}
  {\normalsize\bfseries\sffamily\color{black}}
  {\thesubsection}{0.65em}{#1}
\titlespacing*{\section}{0pt}{11pt}{5pt}
\titlespacing*{\subsection}{0pt}{7pt}{3pt}
\setlength{\parskip}{4.5pt plus 1pt minus 1pt}
\setlength{\parindent}{0pt}
\linespread{1.08}

\usepackage{array}
\usepackage{tabularx}
\usepackage{booktabs}
\usepackage{ragged2e}
\newcolumntype{Y}{>{\RaggedRight\arraybackslash}X}
\newcommand{\code}[1]{\texttt{#1}}
\newcommand{\rx}[1]{\scriptsize\ttfamily #1}

\usepackage{enumitem}
\setlist[itemize]{leftmargin=1.1em,itemsep=1pt,topsep=2pt,parsep=0pt}

\usepackage{listings}
\lstdefinelanguage{ANS}{
  morekeywords={let,function,return,for,in,if,elif,else,print,
                element,compound,reaction,reactants,products,
                formula,name,atomic_mass,valence,symbol,
                check,balanced,energy,exothermic,endothermic,
                molar_mass,compute_yield,mol},
  sensitive=true,
  morecomment=[l]{//},
  morestring=[b]",
  morestring=[b]'
}
\lstdefinestyle{ans}{
  language=ANS,
  basicstyle=\ttfamily\footnotesize,
  keywordstyle=\bfseries\color{anskeyword},
  commentstyle=\itshape\color{anscaption},
  stringstyle=\color{anskeyword!80!black},
  numbers=none,
  frame=none,
  backgroundcolor=\color{anscodebg},
  xleftmargin=5pt,
  xrightmargin=5pt,
  aboveskip=7pt,
  belowskip=7pt,
  showstringspaces=false,
  columns=fullflexible,
  keepspaces=true,
  breaklines=true,
  breakatwhitespace=true,
  upquote=true
}
\lstdefinestyle{ebnf}{
  basicstyle=\ttfamily\scriptsize,
  numbers=none,
  frame=none,
  backgroundcolor=\color{white},
  aboveskip=5pt,
  belowskip=6pt,
  columns=fullflexible,
  keepspaces=true,
  breaklines=true,
  upquote=true,
  literate={->}{{$\rightarrow$}}2
}
\lstset{style=ans}

\usepackage{fancyhdr}
\pagestyle{fancy}
\fancyhf{}
\renewcommand{\headrulewidth}{0pt}
\fancyfoot[L]{\footnotesize\color{anscaption}%
  CSE 341: ANS Design Specification (D1, Part 1) \(\cdot\)
  Student 220104004931 Enes Özkan}
\fancyfoot[R]{\footnotesize\color{anscaption}Page \thepage}

\hypersetup{
  pdftitle  = {ANS Design Specification (D1, Part 1) v1},
  pdfauthor = {220104004931 Enes Özkan},
  pdfsubject= {CSE 341 Project Part 1},
  pdfkeywords={DSL, chemistry, Sebesta, EBNF, lexical analysis}
}

\setcounter{secnumdepth}{0}

\begin{document}

\begin{center}
  {\footnotesize\color{anscaption}CSE 341: Concepts of Programming Languages\par}
  \vspace{3pt}
  {\huge\bfseries\sffamily ANS: Alchemist's Notation Script\par}
  \vspace{4pt}
  {\large Design Specification (D1), Part 1\par}
  \vspace{6pt}
  {\footnotesize\color{anscaption}
    Student: 220104004931 Enes Özkan \(\cdot\)
    Spring 2026 \(\cdot\)
    Gebze Technical University\par}
\end{center}
\vspace{7pt}

\section{4.1\quad Language Overview}

\textbf{Name.} \textbf{ANS}, short for \textit{Alchemist's Notation Script}.

\textbf{Intended domain.} ANS is a small DSL for chemistry programs. Its core objects are
elements, compounds, reactions, and stoichiometric calculations. I designed it so that chemistry
notation stays visible in the program: a user writes \code{H2O}, \code{CH4 + 2*O2}, and
\code{molar\_mass(water)} instead of hiding those ideas inside strings or generic arrays.

\subsection{Sample Program}
\begin{lstlisting}[style=ans]
// Elements and compounds
element H { atomic_mass = 1.008, valence = 1, symbol = "H" };
element O { atomic_mass = 15.999, valence = 2, symbol = "O" };
element C { atomic_mass = 12.011, valence = 4, symbol = "C" };

compound water   { formula = H2O, name = "Water" };
compound methane { formula = CH4, name = "Methane" };
compound co2     { formula = CO2, name = "Carbon dioxide" };
compound oxygen  { formula = O2, name = "Oxygen" };

function molar_mass(compound c) {
    let formula = c.formula;
    let mass = 0.0;
    for (let element in formula) {
        mass += element.atomic_mass * formula[element];
    }
    return mass;
};

reaction combustion {
    reactants = CH4 + 2*O2,
    products  = CO2 + 2*H2O,
    check     = balanced,
    energy    = exothermic 890.3 kJ/mol
};

let mass = molar_mass(water);
if (mass > 18.0) {
    print "Molar mass is above 18.";
} else {
    print "Molar mass is 18 or below.";
}
\end{lstlisting}

\subsection{Sebesta Evaluation Criteria}

Using Sebesta's criteria, ANS mainly chooses \textbf{readability}, \textbf{writability}, and
\textbf{reliability}. Readability matters because the course exam will ask us to defend the design,
and because chemistry code should look close to the equation it represents. Writability matters
because common tasks are short: declaring \code{compound water}, writing \code{2*O2}, and calling
\code{molar\_mass(water)} are all direct operations. Reliability matters most. Formula tokens,
typed parameters, units, and \code{check = balanced} give the compiler places to reject mistakes
before the program runs. I accept the cost: ANS is not general purpose, and the implementation
has to know about formulas, units, and reaction fields. For this project that trade-off is worth it.

The main point I would defend in the exam is that ANS is intentionally narrow. A general
language could calculate the same numbers, but the source code would not show the chemistry
as clearly. In ANS, the syntax itself records which names are elements, which names are
compounds, and which expression is a reaction side.

\section{4.2\quad Lexical Structure}

The scanner is mostly conventional: it recognizes identifiers, literals, keywords, operators,
separators, comments, and whitespace. The chemistry-specific rule is that uppercase tokens are
formula literals, while lowercase names are identifiers or keywords. This keeps \code{H2O} and
\code{water} visually and lexically separate.

{\footnotesize
\begin{tabularx}{\linewidth}{@{}p{2.65cm}p{3.7cm}Y@{}}
\toprule
\textbf{Category} & \textbf{Examples} & \textbf{Pattern or rule} \\
\midrule
Identifiers & \code{water}, \code{mass}, \code{compute\_yield} & \rx{[a-z\_][A-Za-z0-9\_]*} \\
Number literal & \code{20}, \code{1.008}, \code{6.02e23} & \rx{(?:0|[1-9][0-9]*)(?:\textbackslash.[0-9]+)?(?:[eE][+-]?[0-9]+)?} \\
String literal & \code{"Water"} & double quoted text with escapes \\
Char literal & \code{'H'}, \code{'\textbackslash{}n'} & one character or escape in single quotes \\
Formula literal & \code{H}, \code{O2}, \code{H2O}, \code{CH4} & \rx{(?:[A-Z][a-z]?(?:[1-9][0-9]*)?)+} \\
Unit literal & \code{g/mol}, \code{kg/mol}, \code{kJ/mol} & longest match before the slash operator \\
Hard keywords & \code{let}, \code{function}, \code{return}, \code{for}, \code{in}, \code{if}, \code{elif}, \code{else}, \code{print} & identifier text, then keyword table lookup \\
Context words & \code{element}, \code{compound}, \code{reaction}, \code{balanced} & parsed by grammar position \\
Operators & \code{==}, \code{!=}, \code{+=}, \code{=}, \code{+}, \code{-}, \code{*}, \code{/}, \code{<}, \code{>} & longer operators before shorter ones \\
Separators & \code{\{} \code{\}} \code{(} \code{)} \code{[} \code{]} \code{,} \code{;} \code{.} & single character tokens \\
Comments & \code{// comment} & skipped until newline \\
\bottomrule
\end{tabularx}}

Single word units such as \code{mol}, \code{g}, \code{kg}, and \code{kJ} are scanned as identifiers.
The parser treats them as units only after a number. So \code{5.0 mol} is a quantity, but \code{mol}
can still be a normal parameter name in another context. This is simpler than making the lexer
stateful.

This also makes error handling easier to explain. If the scanner sees an uppercase token, it goes
through the formula path. If the parser sees \code{890.3 kJ/mol}, it accepts a quantity because the
unit follows a number. If the same unit appears where an expression cannot contain a quantity,
the parser can report the mistake without relying on hidden lexer state.

\section{4.3\quad Syntax (EBNF)}

This is the Part 1 grammar. It is written for a recursive-descent parser, so precedence and
ambiguous-looking cases are separated into different non-terminals.

\begin{lstlisting}[style=ebnf]
<program> -> { <top_level_statement> } EOF
<top_level_statement> -> <element_declaration> ";"
                       | <compound_declaration> ";"
                       | <reaction_declaration> ";"
                       | <function_declaration> ";"
                       | <variable_declaration> ";"
                       | <assignment_statement> ";"
                       | <augmented_assignment_statement> ";"
                       | <print_statement> ";"
                       | <expression_statement> ";"
                       | <for_statement> [ ";" ]
                       | <if_statement>

<block> -> "{" { <statement> } "}"
<statement> -> <variable_declaration> ";"
             | <assignment_statement> ";"
             | <augmented_assignment_statement> ";"
             | <return_statement> ";"
             | <print_statement> ";"
             | <expression_statement> ";"
             | <for_statement>
             | <if_statement>

<variable_declaration> -> "let" <identifier> "=" <expression>
<assignment_statement> -> <assignable> "=" <expression>
<augmented_assignment_statement> -> <assignable> <augmented_assignment_operator> <expression>
<augmented_assignment_operator> -> "+=" | "-=" | "*=" | "/="
<return_statement> -> "return" [ <expression> ]
<print_statement> -> "print" <expression>
<expression_statement> -> <expression>

<function_declaration> -> "function" <identifier> "(" [ <parameter_list> ] ")" <block>
<parameter_list> -> <parameter> { "," <parameter> }
<parameter> -> [ <type_name> ] <identifier>
<type_name> -> "element" | "compound" | "reaction"

<for_statement> -> "for" "(" "let" <identifier> "in" <expression> ")" <block>
<if_statement> -> "if" "(" <expression> ")" <block> { <elif_clause> } [ <else_clause> ]
<elif_clause> -> "elif" "(" <expression> ")" <block>
<else_clause> -> "else" <block>

<element_declaration> -> "element" <element_symbol> "{" <element_field_list> "}"
<element_symbol> -> <formula_literal>
<element_field_list> -> <element_field> { "," <element_field> } [ "," ]
<element_field> -> "atomic_mass" "=" <number_literal>
                 | "valence" "=" <number_literal>
                 | "symbol" "=" <string_literal>

<compound_declaration> -> "compound" <identifier> "{" <compound_field_list> "}"
<compound_field_list> -> <compound_field> { "," <compound_field> } [ "," ]
<compound_field> -> "formula" "=" <formula_literal>
                  | "name" "=" <string_literal>

<reaction_declaration> -> "reaction" <identifier> "{" <reaction_field_list> "}"
<reaction_field_list> -> <reaction_field> { "," <reaction_field> } [ "," ]
<reaction_field> -> "reactants" "=" <reaction_expression>
                  | "products" "=" <reaction_expression>
                  | "check" "=" <balance_check>
                  | "energy" "=" <energy_expression>
<balance_check> -> "balanced"
<energy_expression> -> <energy_kind> <quantity_literal>
<energy_kind> -> "exothermic" | "endothermic"

<reaction_expression> -> <reaction_term> { "+" <reaction_term> }
<reaction_term> -> <species_ref> | <coefficient> "*" <species_ref>
<coefficient> -> <number_literal>
<species_ref> -> <formula_literal> | <identifier>

<expression> -> <comparison_expression>
<comparison_expression> -> <additive_expression>
                           [ ( "==" | "!=" | ">=" | "<=" | ">" | "<" ) <additive_expression> ]
<additive_expression> -> <multiplicative_expression>
                         { ( "+" | "-" ) <multiplicative_expression> }
<multiplicative_expression> -> <unary_expression>
                               { ( "*" | "/" ) <unary_expression> }
<unary_expression> -> [ "-" ] <postfix_expression>
<postfix_expression> -> <primary_expression> { <postfix_operator> }
<postfix_operator> -> "." <identifier> | "[" <expression> "]" | "(" [ <argument_list> ] ")"
<primary_expression> -> <number_or_quantity_literal>
                      | <string_literal>
                      | <char_literal>
                      | <list_literal>
                      | <formula_literal>
                      | <identifier>
                      | "(" <expression> ")"

<argument_list> -> <expression> { "," <expression> }
<number_or_quantity_literal> -> <number_literal> [ <unit> ]
<quantity_literal> -> <number_literal> <unit>
<unit> -> UNIT_LITERAL | "mol" | "g" | "kg" | "kJ"
<list_literal> -> "[" [ <argument_list> ] "]"

<assignable> -> <identifier> { <assignable_postfix> }
<assignable_postfix> -> "." <identifier> | "[" <expression> "]"

<identifier> -> IDENTIFIER
<number_literal> -> NUMBER_LITERAL
<string_literal> -> STRING_LITERAL
<char_literal> -> CHAR_LITERAL
<formula_literal> -> FORMULA_LITERAL
\end{lstlisting}

\subsection{Ambiguity Notes}
\begin{itemize}
  \item Arithmetic precedence is encoded by grammar levels, not by a later fix-up step.
  \item Reaction expressions use their own grammar. In \code{CH4 + 2*O2}, the \code{*} belongs to the coefficient \code{2}.
  \item ANS requires \code{2*O2}, not \code{2O2}, so coefficients and formula subscripts are not mixed.
  \item Formula literals start uppercase and identifiers start lowercase or underscore.
  \item \code{element}, \code{compound}, and \code{reaction} become type names only inside parameter lists.
  \item Braced blocks remove the usual dangling \code{else} problem.
\end{itemize}

The grammar also separates assignment targets from general expressions. Function calls are
valid expressions, but they are not valid assignment targets. That keeps a line such as
\code{f(x) = 3} invalid instead of forcing the language to define assignment to a function result.

\section{4.6\quad Names, Binding, Scope, and Lifetime}

\textbf{Identifiers.} A legal identifier matches \code{[a-z\_][A-Za-z0-9\_]*}. ANS is
case-sensitive. This is also how the language separates names from chemistry notation:
\code{water} is an identifier, but \code{H2O} and \code{CO2} are formula literals. The hard
reserved words are \code{let}, \code{function}, \code{return}, \code{for}, \code{in}, \code{if},
\code{elif}, \code{else}, and \code{print}. Chemistry words such as \code{compound} and
\code{reaction} are contextual.

\textbf{Binding times.} Most reliability-related bindings happen before execution. The lexer
binds lexemes to token categories. The parser binds names to declaration forms and parameters
to optional type annotations. The later semantic phase binds compounds and reactions to their
declarations, validates formulas, and checks \code{check = balanced}. Runtime binding is still used
for argument values, local storage, loop iteration values, and ordinary variable values.

{\footnotesize
\begin{tabularx}{\linewidth}{@{}p{4.1cm}p{3.2cm}Y@{}}
\toprule
\textbf{Binding} & \textbf{When} & \textbf{Example} \\
\midrule
Lexeme to token category & Lex time & \code{water} is an identifier, \code{H2O} is a formula. \\
Identifier to declaration & Compile time & \code{mass} resolves by static scope. \\
Parameter to declared type & Compile time & \code{compound c} marks \code{c} as a compound parameter. \\
Reaction balance decision & Compile time & \code{check = balanced} asks for atom-count validation. \\
Parameter to argument value & Run time & \code{molar\_mass(water)} binds \code{c} to \code{water}. \\
Local variable to storage & Run time & \code{let mass = 0.0} allocates during the call. \\
\bottomrule
\end{tabularx}}

\textbf{Scope.} ANS uses static scoping. A name is resolved by looking in the current block and
then outward in the source text. I chose this because chemistry calculations should not depend
on who called a function. If lookup were dynamic, a local name in the caller could silently change
what \code{molar\_mass} or \code{compute\_yield} means.

For example, if a caller had a local variable named \code{water}, dynamic scoping could make a
function read that caller's binding instead of the compound declaration the function was written
against. That would be confusing in any language, but in a chemistry DSL it is especially bad
because the same expression could refer to a different compound depending on the call chain.

{\footnotesize
\begin{tabularx}{\linewidth}{@{}p{4.2cm}p{3.1cm}Y@{}}
\toprule
\textbf{Entity} & \textbf{Lifetime} & \textbf{Reason} \\
\midrule
Top-level elements, compounds, reactions & Static & They describe the program's chemistry model. \\
Top-level \code{let} variables & Static & They live for the whole run. \\
Function parameters & Stack-dynamic & Each call gets its own arguments. \\
Local \code{let} variables & Stack-dynamic & They disappear when the block exits. \\
\code{for} loop variable & Stack-dynamic & It belongs only to the loop body. \\
Lists and chemistry objects & Implicit heap-dynamic & The runtime manages them. ANS has no \code{new}, \code{delete}, or pointers. \\
\bottomrule
\end{tabularx}}

This lifetime design follows the reliability goal. The programmer works with formulas and
reactions, not manual memory management.

\end{document}
